\cotefile{57}
\scanpage{193}
\uline{Remarques}. 1) On notera que la démonstration
montre que si $\dim X=1$, alors $\mathrm{Pic}(\widehat E)\unsure{-|\to}\mathrm{Pic}(\deleted{\unsure{E|\widehat E}}X)$
a un noyau qui est connexe\inserted{ et $\unsure{=|\simeq}$ unipotent}, \deleted{\unsure{le|la}} donc
$\mathrm{Pic}(S-\unsure{\bullet|o|a})\unsure{-|\to}\mathrm{Pic}(\widehat S-\bullet)$ induit un isom sur les
groupes \uline{mod composantes connexes} de l'identité
[qui dans ce cas est fini, isom. à $\mathbb Z/d\mathbb Z$ où
$d$ est le degré projectif de $X$]\inserted{ et $=$ modulo partie unipotente du \unsure{dernier.|2\textsuperscript{d} terme.}} Le $=$ résultat
semble valable si $X$ est de dim quelconque,
lorsque \deleted{\unsure{on se met|on se}} la caractéristique \deleted{résiduelle}
$p$ de $\deleted{\unsure{k}}k$ est $>0$, et qu'on raisonne
\uline{mod} \uline{$p$-groupes}\inserted{ \uline{de $p$-torsion}}; i.e. le groupe de N.S. local
semble le $=$ pour $S$ et $\widehat S$, modulo des
$p$-groupes [ext. d'un $p$-groupe fini par un
groupe alg. unipotent sur $k$]; \unsure{utiliser|utilise|utilisant} le
fait que les obstructions dans des
% slanted
\margin{Regarder aussi le cas de la car. résiduelle nulle \ldots\,.}
$H^2(X,\underline O_X(n))$ qu'on rencontre sont nulles
pour $n$ grand, et sont annulés par $p$ \ldots\,. 
Ceci est en accord avec la théorie des
groupes \unsure{fondamentx|fondamentaux} \ldots\,. 
% archive sheet number: 192

\scanpage{194}
2) On peut aussi regarder le \unsure{hensélisé|hensilisé}
$\widetilde S$ de $S$, et les hom
\[
\mathrm{Pic}(S-a)\longrightarrow\mathrm{Pic}(\widetilde S-\widetilde a)-\mathrm{Pic}(\widehat S-\widehat a).
\]
Il est possible que la dernière flèche
écrite soit toujours un isomorphisme.
On le voit immédiatement, en tous cas,
lorsque $X$ est de dim $1$ (\unsure{i.e.|càd} $S$ de
dim $2$), par le raisonnement habituel
sur l'éclaté \ldots\,. \unsure{Donc|Dans} l'exemple précédent
des courbes planes donne autant
d'exemples où \unsure{les|le} \unsure{gpes|groupes|groupe} de Picard local de
$S$ et celui de son hensélisé ne sont
pas isomorphes \ldots\,.
% archive sheet number: 193

\scanpage{195}
$G^\tau$, $G^0$ etc

Structure \inserted{locale} du $\underline{\mathrm{Pic}}_{X/S}$

\deleted{(\unsure{finitude}, \unsure{simplicité|simplification})}
% archive sheet number: 194

\scanpage{196}
\fbox{\shortstack{Continuité de la\\Torsion}}

\diagram{\begin{tikzcd}[ampersand replacement=\&,row sep=small]G\arrow[d,dash]\\ S\end{tikzcd}}
schéma en groupes \inserted{comm.} projectif sur $S$ \deleted{\unsure{loc.|loc}} noethérien \uline{local}
\deleted{$G^0$ univ. ouvert sur $S$.}
$s\in S$ \inserted{pt fermé}, $s_i$ générisations maximales.
Pour tt $t\in S$, soit $G_t^0$ muni de la structure \uline{abélienne} induite, soit $T_t(G)=G_t/G_t^0$
et $\sigma(t)=\mathrm{rang}_{k(t)}T_t(G)$

Implications

(i) \deleted{$G$ \uline{plat} sur $S$,} $G^{\deleted{\unsure{0|\circ}}}$ \deleted{contient}\inserted{admet} un sous-schéma \uline{abélien} $A$ ayant
$G^0$ comme sous-ens. sous-jacent, et $G$ plat sur $S$

\margin{$\Updownarrow$}(ii) $G$ admet un $A$ comme dessus, et $T(G)\simeq G/A$ est
\uline{plat} sur $S$ [N.B. il est déjà fini sur $S$]

\margin{$\Downarrow$}(iii) Pour tout $t\in S$, on a \deleted{rang} $\sigma(s)\deleted{\unsure{\geq}}=\sigma(\deleted{\unsure{s}}t)$ \deleted{[\unsure{ou seul.} $\sigma(s)\leq\sigma(s_i)$]}

\margin{$\Updownarrow$}(iv) Pour tt \deleted{\unsure{$t\in S$, on}} \unsure{$i$, on a} $\sigma(\unsure{s|t})=\sigma(s_i)$ [ou seulemt $\sigma(s)\leq\sigma(s_i)$]

Si $G^0$ univ. ouvert sur $S$, \inserted{i.e. les fibres de $G$ ont ttes $=$ dim.} on a \textup{(iii)}$\Longleftrightarrow$\textup{(iv)}.

Si \inserted{de plus} $S$ est \uline{réduit}, \uuline{ttes} les conditions envisagées sont équivalentes.

On sait que (\unsure{théorie|Théorie} des \unsure{passages|groupes} au quotient \inserted{plat}) que si $A$
existe, alors $G\longrightarrow G/A$ est fid. plat et projectif, et que
alors $G$ plat sur $S$ sss $G/A$ plat sur $S$, ce qui donne
l'équivalence de (i) et (ii). De plus, si $A$ existe, alors
on a pour tt $t\in S$, $T_t(G)\simeq T(G)_t$ où $T(G)=G/A$.
Donc (ii) $\Longrightarrow$ (iv), \unsure{à|:|a} fortiori (ii) $\Longrightarrow$ (iii).

Pour (iii) $\Longleftrightarrow$ (iv), on \deleted{\unsure{aura besoin du}}\inserted{applique le}

\uline{Lemme} \inserted{Supposons $G^0$ univ. ouvert sur $S$, i.e. les fibres de $G$ ont ttes $=$ dim.} La fonction $\deleted{\unsure{f}}\sigma$ est \uline{semi-continue} \uline{supérieurement}.
Comme elle est évidemment constructible (cf l'argument
précédent \ldots\,.) on est ramené à prouver, dans le cas où $S$ est
le spectre d'un anneau de val. discrète, $s'$ son pt gén.,
que $\sigma(s)\geq\sigma(s')$. Soit $G^0$ muni de la structure adhérence
de celle de $G^0_{s'}$, donc $G^0$ est un \uline{ss-groupe} de $G$ \uline{plat}$/S$,
et on peut former le groupe quotient $Q=G/G_0$, qui est
fini sur $S$. Alors
\[
\sigma(s)\geq\mathrm{rang}\,[Q_s:k(s)]\geq\mathrm{rang}\,[Q_{s'}:k(\unsure{s|s'})]=\sigma(s')
\]
d'où l'inégalité voulue. Noter que la première inégalité est
une égalité sss $G_s^0$ est réduit, i.e. $G^0$ est \uline{abélien} sur $S$, et
que la deuxième l'est sss $Q$ est plat sur $S$. / Cela prouve ainsi
que dans le cas $S=\mathrm{Spec}(V)$, on a (iv) $\Longrightarrow$ (i). C'est aussi
vrai dès que $S$ est seulement réduit, En effet, on sait cons%
% margin note written sideways
\margin{truire comme sous-schéma du schéma de Hilbert de $G/S$, un schéma $S'\to S$ résolvant le Pb universel correspondant à la recherche d'un $A$. Alors $S'\unsure{-|\to}S$ est un mono, et surjectif, prouvons que c'est un isom, i.e. \uline{propre}. Le critère valuatif nous ramène \unsure{au cas|encore à} $S=\mathrm{Spec}(V)$, où les remarques précédentes s'appliquent. D'où résulte l'existence du $A$, donc du $T=G/A$. Enfin $T$ est plat sur $S$ en vertu de (iv), [\ldots{} critère \unsure{val.|util.} \ldots].}
% archive sheet number: 195

\scanpage{197}
Nous allons donner quelques \unsure{rappel\deleted{a}s|rappels} d'Algebre Commutative:

(i) Soit $A$ un anneau \uline{local} \uline{noetherien}. Alors $\dim\,(A)={}$ nombre \unsure{minimmum|minimum}
de generateurs d'un idea\deleted{\unsure{m}}l de definition de $A$. En particulier $\dim(A)$
est finie. Si $A$ \inserted{donc} est un anneau noetherien \unsure{quelconaue|quelconque}, $\underline p$ un ideal premier
de $A$, comme le rang de $\underline p$ est egal evidemment a $\dim A_{\underline p}$, il
s'ensuit qu'il est fini.

(ii) Si $A$ est un anneau \deleted{\unsure{regulier}}\inserted{local} regulier, il satisfait la condition
des chaines (on peut referer a l'article de Serre au Japon). C\uline{'}est donc
vrai aussi par suite pour un anneau regulier quelconque (pas nec.
local), donc pour un quotient d'un tel anneau, i.e. pratiquement pour
\unsure{\deleted{b}tous|bous} les anneaux noetheriens rencontres dans la nature.

(iii) Soient $k$ un corps, $A$ une algebre integre de type fini sur $k$,
$m$ un ideal maximal de $A$, alors $\dim A_m={}$ degre de \unsure{transcendance|transcsndance} de $K$ sur
$k$, ou $K$ est le corps des fractions de $A$. Par suite $A$ est \unsure{equidimansionnel|equidimensionnel}
(et meme strictement, en vertu de (ii)) de \unsure{dimsnsion|dimension} egale au dit
degre de transcendance.

(iv) Enfin, nous utiliserons au Chap.III la definition de la dimension
d'un anneau local noeth. $A$ par le \inserted{degre du} polynome de Hilbert-Samuel,
qu'il\deleted{x} faut donc rappeler ici.

En langage geometrique, (i) signifie que si on a un preschema
\deleted{\unsure{irreductible $X$ et|irreductible $X$}} noetherien $X$ et une partie fermee irreductible $Y$,
alors si $X$ est affine, $\mathrm{codim}_X Y$ est egal au nombre minimum de sections
de $\underline O_X$ telles que $Y$ soit une composante de l'ensemble des zeros communs
a ces sections; et en tous cas ($X$ affine ou non) cette codimension est
finie. (ii) signifie que pratiquement, les preschemas noetheriens
qu'on rencontre satisfont la condition des chaines. Quant a (iii), il
implique qu'un \unsure{presbhama|preschema} algebrique irreductible est strictement
equidimensionnel, et que sa dimension est identique a celle de
toute partie ouverte non vide de $X$ (ce qui ne serait pas vrai par
exemple de $\mathrm{Spec}(V)$, ou $V$ est un anneau de valuation discrete, bien
que $\mathrm{Spec}(V)$ soit manifestement strictement equidimensionnel). \unsure{Enonçnons|Enonçons}

\uline{Proposition} 7. Soit $X$ un preschema algebrique sur le corps $k$, supposons
$X$ irreductible de point generique $z$. Alors $X$ est strictement equidimensionnel
et de dimension egal a $\mathrm{deg.tr.}_k k(z)$ .

\uline{Corollaire} 1 \quad La dim. de $X$ est egale a celle de toute partie ouverte
non vide. \qquad - \quad On en deduit:

\uline{Corollaire} 2 \quad Soit $X$ un \unsure{preschama|preschema} algebrique (pas nec. irred.). Pour
tout point \deleted{ferme} $x$ de $X$, on a $\dim_x X=\mathrm{Sup}\,\dim\,(X_1)$ ou $X_1$ parcourt les
composantes de $X$ contenant $x$. Si $x$ est ferme, on a $\dim_x X=\mathrm{codim}_X x$
% archive sheet number: 196

\scanpage{198}
\margin{$?\,\left\{\vphantom{\begin{matrix}a\\a\\a\end{matrix}}\right.$}
\uline{Corollaire} (\deleted{\unsure{Soit}}les fibres de $G$ ont $=$ dimension). Les conditions
(iii) et (iv) équivalent aussi à la suivante, où
$n$ est un entier fixé \deleted{\unsure{tt|tel que}} multiple des $\sigma(t)^2$; et
$\,{}_nG=\mathrm{Ker}(G\overset n{\unsure{-|\longrightarrow}}G)$
\[
\operatorname{rang}_{k(s)}{}_nG_s=\operatorname{rang}_{k(s_i)}{}_nG_{s_i}\qquad\text{pour tt }i,\tag*{\textup{(\unsure{v|w})}\quad}
\]
ou
\[
\operatorname{rang}_{k(s)}{}_nG_s=\operatorname{rang}_{k(t)}{}_nG_t\qquad\text{pour tt }t\in S.\tag*{\textup{(v bis)}\quad}
\]

\[
H^1(\mathbb W_{\unsure{nn|nm}}\qquad\qquad \unsure{F^n|f^n}V^n=0
\]

\diagram{\begingroup\scriptsize\setlength{\arraycolsep}{1.2pt}\(\begin{array}{ccc}
\hline
\multicolumn{1}{|c}{\operatorname{long} H^1(\overline X_s,\mathbb Z/p^n)}&
{}+\operatorname{long} H^1(\overline X_s,\mu_{p^n})&
\multicolumn{1}{c|}{{}+\deleted{\unsure{\dim}}\operatorname{long}_{\mathbb W_n^{\unsure{\sigma|o}}[F,V]}H^1(\overline X_s,\mathbb W_{n,n})}\\
\multicolumn{1}{|c}{\text{partie multiplicative}}&\text{partie additive}&\multicolumn{1}{c|}{/}\\
\multicolumn{1}{|c}{}&\text{\unsure{séparable|unipotente}}&\multicolumn{1}{c|}{}\\ \hline
\Vert&\Vert&\Vert\\
\mathrm{Hom}(\pi_1(\overline X_s),\mathbb Z/p^n)&{}_{p^n}\mathrm{Pic}(\overline X_s)&\mathrm{Ker}\bigl(H^1(\overline X_s,\mathbb W_n)\xrightarrow{F^n} H^1(\unsure{\overline X_s|X_s},\mathbb W_n)\bigr)
\end{array}\)\endgroup}
% archive sheet number: 197

\scanpage{199}
Meme pour un schema \inserted{local} noetherien, il est possible que la condition
des chaines ne soit pas verifiee, cf Nagata, On the chain problem of
prime ideals, Nagoya Math. Journal, Vol. 10, June 1956, \deleted{\unsure{qui contient|qui traitant}}
(article contenant des resultats interessants sur cette question).

Voir aussi

Nous \deleted{\unsure{verrons}}\inserted{rappelerons} cependant au numero suivant des conditions tres generales
moyennant lesquelles la condition des chaines est satisfaite, conditions
heureusement tou\deleted{\unsure{g}}jours satisfaites en pratique. Quant a l'equidimensionnalite
stricte, elle peut etre en defaut dans des cas non
pathologiques (donc lors meme que la condition des chaines est satisfaite)
pour deux raisons: 1) Les composantes irreductibles de $X$ n'ont
pas meme dimension \deleted{\unsure{et}} (\unsure{on|en} dit alors que $X$ n'est pas \uline{equidimensionnel}\deleted{\unsure{)}},
terminologie qu'il faudrait introduire dans le texte -) 2) Meme si
$X$ est irreductible, et meme si $X$ est un schema integre \uline{regulier}, (i.e.
dont les anneaux locaux sont des anneaux reguliers), il est possible
que les points fermes de $X$ n'aient pas meme codimension. Ainsi soit $V$
un anneau de valuation discrete, soit $K$ son corps des fractions, \inserted{$k$ son corps residuel} $T$ une
indeterminee,posons $A=V\ T$, $X=\mathrm{Spec}(A)$, de sorte que les points
fermes de $X$ correspondent aux ideaux \unsure{maxmaux|maximaux} de $A$, et leur codimension
dans $X$ au \uline{rang} desdits ideaux. Il y a des ideaux maximaux de \deleted{rang}
rang 2, par exemple ceux provenant d'ideaux maximaux de $k\ T$, \deleted{\unsure{ou de}}
\deleted{\uline{certains} ideaux maximaux de \unsure{$K\ T$|$k\ T$}.} Mais \deleted{\unsure{pour}} il y a des ideaux
\unsure{makkmaux|maximaux} de $A$ de rang 1, provenant de\deleted{\unsure{s autres|s}} \inserted{certains} ideaux \unsure{makimaux|maximaux} de $K\ T$,
par exemple l'ideal principal $(aT-1)$, ou $a$ est un generateur de
l'ideal maximal de $V$.

2. \uline{Dimension des preschemas.}

Rappelons la definition suivante:

\uline{Definition} 2. Soient $A$ un anneau commutatif. On appelle dimension de
$A$ la dim. de $\mathrm{Spec}(A)$, et si $J$ est un ideal de $A$, on appelle rang de $J$
la codimension de $V(J)$ dans $\mathrm{Spec}\unsure{!|(}A)$. On dit que $A$ \deleted{est} satisfait
la condition des chianes (resp. est equidim. resp. strictement equidim)
si c'est comme ça pour $\mathrm{Spec}(A)$.

\uline{Proposition 6} \quad Soient $X$ un preschema, $Y$ une partie fermee irreductible
alors on a
\[
\mathrm{codim}_X Y=\dim\underline O_{X/Y}
\]

\uline{Corollaire} 1 \quad La condition des chaines pour $X$ equivaut a la condition
des chaines pour les $\underline O_x$ ($x\unsure{\epsilon|\in}X$); si $X$ est noetherien \deleted{\unsure{on|$X$}}, on peut se
borner aux point fermes.

\uline{Corollaire} 2 \quad Si $X$ est noetherien, on a $\dim X=\mathop{\mathrm{Sup}}\limits_{x\;X}\dim\underline O_x$, et
on peut s'y borner aux point fermes.
% archive sheet number: 198

\scanpage{201}
\fbox{Sous-groupe $G^\tau$}

Soit $G$ un schéma en groupes sur le \uline{corps} $k$,
localement de type fini sur $k$. Alors :

\deleted{1°) La composante connexe de l'élément neutre $e$
est de type fini sur $k$, et \unsure{relt.|ablt.} irréductible}

\deleted{\unsure{2}}1°) Les composantes connexes \inserted{$G_i$} de $G$ sont de type
fini sur $k$, et irréductibles

2°) La composante connexe \inserted{$G^{\deleted{\unsure{0}}}_0$} de $e$ est
\unsure{relt.|ablt.|géom.} irréductible, et c'est un sous-groupe invariant

3°) \deleted{Si $k$ est alg. clos, alors les \unsure{$G_i$} \unsure{ouvert et fermé} les composantes
connexes de $G$ sont les parties de la forme
$aG^{\unsure{0|\circ}}_0$, où \unsure{$a$} $a\in G(k)$} \inserted{Les $G_i$ sont stables sous $G_0$. Si $G_i$ contient un pt $a$ rat. sur $k$, alors $G_i=G_0a$}

4°) Le \deleted{groupe}\inserted{schéma quotient} $G/G^{\deleted{\unsure{0|\circ}}}_0$ existe il est discret i.e.
localement fini sur $k$, $G\unsure{-|\to}G/G^{\deleted{\unsure{0|\circ}}}_0$ est
un morphisme \deleted{\unsure{plat}} fid plat. $G/G_0$ est
un $k$ groupe, sépa\uuline{rable} et loc. fini. \deleted{\unsure{sur $k$.}}

\deleted{La composante connexe $G^0$ de $e_{\bar k}$ dans $G_{\bar k}$ est}
Supposons d'abord $k$ alg. clos. \deleted{Soit $G_0$}
\deleted{définie sur} Comme $G_0$ est connexe,
il en est de $=$ de $G_0\times G_0$ (\uuline{référence} \ldots)
% slanted
\margin{$\times$ Attention, $G_0$ n'est pas à \uline{priori} de type fini.}
donc $G_0\times G_{\unsure{0|e}}\unsure{-|\to}G$ induit $G_0\times G_0-G_0$,
i.e. $G_0$ est un sous-groupe. Soit $U$
un voisinage ouvert \inserted{\unsure{affine}} de $e$, je dis que
$U^2=G_0$ \deleted{\unsure{et}} donc $G_0$ est \unsure{quasi|qu.} cpct et
% archive sheet number: 200

\scanpage{202}
par suite de type fini sur $k$. Soit en effet
$x$ un pt de $G_0$ rat sur $k$, il suffit de
prouver que $x\in U^{\unsure{2|?}}$ [N.B. $U^2$ est constructible
qu. cpct dans $G_0$] i.e. \deleted{$xU\cap U=$} l'existence
de $y$ et $z\in U(k)$ tels que $yz=x$, i.e.
$y=xz^{-1}$; ce qui signifie que $(U\cap xU)(k)\ne\varnothing$,
or $U$ et $xU$ sont des ouverts non vides
dans $G_0$, \deleted{\unsure{donc}}. Donc $G_0$ est de type
fini sur $k$. S'il n'était irréductible, \uline{tt}
point fermé de $G_0$ appartiendrait à \deleted{une
composante} deux composantes irréductibles
distinctes (par translation) absurde. On voit
aussi \deleted{\unsure{de suite}} 3° \deleted{\unsure{que}}\unsure{par|pour} \unsure{transp\textsuperscript{t}|transpos|transport} de
\deleted{Considérons la projection $G_k$}
structure [il suffit de considérer les
composantes connexes des pts fermés]
d'où aussi 1°. \deleted{\unsure{On obtient ainsi}} \inserted{\deleted{\unsure{Cela}} Prenant $U=G_0$, on}
voit que $G_0$ est un \uuline{sous-groupe}. Il est invariant
car \ldots

\deleted{Pour} Soit $k$ quelconque. \deleted{Soit} Il
suffit encore d'étudier les composantes
connexes \inserted{$U$} des pts \uline{fermés} $x$. \deleted{\unsure{Un tel}
pt devient rationnel sur une extension
finie $k'$, donc la composante connexe
\unsure{se décompose|décompose} en un nb fini}
% archive sheet number: 201

\scanpage{203}
Il n'y a qu'un nb fini de pts de
$G_{\bar k}$ au dessus, \deleted{soit $U'$} la réunion des
de leurs composantes connexes est $U_{\bar k}$ [\uline{réf}]
\margin{$\times$}\deleted{donc $U$ \unsure{est q}} qui est de type fini sur
$\bar k$, donc $U$ l'est sur $k$. Comme $U_{\bar k}$ est
loc. irréductible, il en est de $=$ de $U$
\margin{$\times$}[réf]. Comme $e$ est rat/$k$, il s'ensuit
que sa composante connexe \inserted{\uline{\unsure{absolue}}} est définie sur $k$
[réf] donc \unsure{relt.|ablt.} irréductible; \deleted{et on voit} \inserted{Comme}
$G_{0\,\bar k}$ est invariant, $G_0$ est invariant.

\noindent\rule{1.15cm}{0.3pt}

Pour prouver l'existence du préschéma
quotient, on note plus généralement

\uline{Corollaire} \quad $G$ schéma en groupes loc de
type fini sur $k$, \deleted{\unsure{soit|G}}$H$ sous-schéma en
groupes ouvert \deleted{invariant} alors $H$ est
fermé, \deleted{Si $H$ est invariant, alors} $G/H$
existe, et $G\unsure{-|\to}G/H$ est fid. plat. (c'est
donc un \uuline{groupe} si $H$ est distingué)
% slanted
\margin{se généralise : \unsure{sur|une} base plus \unsure{général|générale}, quand $\deleted{\unsure{G}}H$ est \uuline{plat} sur \unsure{$k$|$S$} ? Oui si $S$ est \uuline{artinien}.}

\deleted{On peut utiliser une construction \unsure{spéciale} dans
le cas $H=G_0$.} Dans le cas $H=G_0$
on prend une composante connexe de $G$,
elle est stable sous $G_0$ [à signaler \ldots]
et \deleted{\unsure{puis|qui}} $G/G_0$ existe \deleted{\unsure{génériqut|géométriqut|groupe}}, car on le
voit trivialemt après l'extension finie de $k$.
% archive sheet number: 202

\scanpage{204}
Le cas $H$ quelconque s'en déduit

\deleted{Soit} Autres énoncés : $G$ simple sur $k$
$\Longleftrightarrow$ $G$ séparable sur $k$ \unsure{;|,} $\Longleftrightarrow$ $G$ séparable sur $k$ en $e$

\begin{center}$-\;-\;-\;-\;-\;-\;-$\\[-2pt]\rule{2.5cm}{0.3pt}\end{center}

Supposons $G$ \uline{commutatif}, \deleted{Soit} alors
% slanted
\margin{inutile cf plus bas}
$\Gamma=G/G^0$ est un \deleted{\unsure{préschéma}}\inserted{\uuline{discret}} $k$-groupe commutatif \inserted{séparable sur $k$}
donc $\Gamma\xrightarrow{\times n}\Gamma$ est un hom de groupe,
dont le noyau $\,{}_n\Gamma$ est un sous-groupe
fermé.

\uline{Proposition} \deleted{Il existe un plus grand entier} \inserted{Soit $\Gamma$ un $k$-\unsure{groupe|gpe}, loc. fini sur $k$. La suite des $\,{}_n\Gamma$ est}
localement constante, et \deleted{\unsure{c'est un sous}}\inserted{son sup,,}
soit $\deleted{\unsure{a}}\Gamma^\tau$, est un sous-groupe ouvert de
$\Gamma$. Pour que \deleted{$a\in$} l'élément $a$ de $\Gamma$
soit dans $\Gamma^\tau$, il faut. \unsure{so.|su.} qu'il
existe un $n$ tel que $na=e$.
% slanted
\margin{après \unsure{un|une} \unsure{suite des|mot des|mot sur les} groupes \uline{finis} et des sous-groupes $\Gamma$}

Soit $a\in\Gamma$, \deleted{tel que} si $a\in{}_n\Gamma$ on a
$na=0$. Réciproqt, supposons $na=0$,
comme il existe $m$ tel que $m\Gamma_{a\deleted{\unsure{k}}}=0$
[réf] il s'ensuit que \deleted{\unsure{mor.}} $\{a\}\deleted{\unsure{u}}\xrightarrow{\mathrm{incl.}}\Gamma\xrightarrow{mn}\Gamma$
% archive sheet number: 203

\scanpage{205}
est nul, où \deleted{\unsure{nous}}\inserted{$\{a\}$} est le schéma ouvert
induit sur $a$, donc $\{a\}$ est dominé
par $\,{}_{(mn)}\Gamma$. Cela prouve bien que
la suite des $\,{}_n\Gamma$ est loc. const, que son
sup est un ouvert, $\Gamma^\tau$ qui est caractérisé
comme on l'a dit. Que cet ouvert
soit \deleted{invariant}\inserted{un sous-groupe} est immédiat, et résulte aussi du

\deleted{N.B.} \uline{Corollaire}. Pour tout $S$ \deleted{\unsure{local}}\inserted{qu. cpct} sur $k$,
$\Gamma^\tau(S)={}$ torsion de $\Gamma(S)$. \deleted{\unsure{Et les fibres|Et la plus}}

\deleted{\uline{Corollaire}. $\Gamma^\tau$ est le plus petit sous-groupe
fermé de $\Gamma$ tel que $\Gamma/\Gamma^\tau$ soit séparable sans \unsure{tor.}}

Si $G$ est un \deleted{\unsure{sché}} $k$-groupe
\deleted{quelconq} loc de type fini, on
désigne par $G^\tau$ l'image inverse de
$\Gamma^\tau$ par $G\unsure{-|\to}G/G_0=\Gamma$.

\uline{Proposition} $G^{\deleted{\unsure{0|\tau}}}_0$ \inserted{(resp. $G_{\unsure{\tau|\varepsilon}}^{\deleted{\unsure{\tau|0}}}$)} est le plus petit sous-groupe
fermé de $G$ tel que $G/G_0$ soit
\uline{discret} \uline{séparable} [resp. discret \deleted{séparable}\inserted{sans}
torsion]. C'est aussi le sup. \deleted{\unsure{$\bigcup_n U_n$}}\inserted{\unsure{(croissant)|croissant)} des sous-groupes [ouverts; \unsure{ou : ouvert} \unsure{``\quad''} }]
de $G$ \unsure{(de} \deleted{\unsure{type}} \unsure{finis,} \unsure{resp. les}
% slanted
\margin{à prouver \ldots}
$G$ est dit \uline{``sans torsion''} si
$(G/G_0)^\tau=\{e\}$, \deleted{\unsure{Alors}} $G^\tau$ est dit le
\uline{sous-groupe de torsion de $G$}
% archive sheet number: 204

\scanpage{206}
$\deleted{\unsure{\Gamma}}G$ \quad $k$-gpe loc. de type fini sur $k$,
$G_0$ \quad composante connexe
\[
\Gamma=G/G_0
\]
\[
\Gamma\xrightarrow{\pi_n}\Gamma\qquad {}_n\Gamma=\pi_n^{-1}(e)
\]
c'est un sous-schéma \uline{ouvert} et \uline{fermé} [car $\Gamma$
discret \uline{sép}] ens des $a$ tels que $\pi_n(a)=e$

(i) La suite des $\,{}_n\Gamma$ est loc. stationnaire, \deleted{\unsure{or}} (\unsure{trivial}
car $\Gamma$ est discret] \deleted{\unsure{alors}} son
sup est un sous-schéma \uline{ouvert} \uline{induit} $\Gamma^\tau$
de $\Gamma$ [N.B. c'est vrai pour tt $k$-groupe
loc. fini sur $k$, pas néc. séparable, mais
c'est moins évident, le cas radiciel
demande une démonstration spéciale]

\deleted{\unsure{(ii)}} Soit $a\in\Gamma$, alors $a\in{}_n\Gamma\deleted{\unsure{^\tau}}$ sss \deleted{\unsure{$a=e$}}
\deleted{il existe $n$}\inserted{$\pi_n(a)=e$}
$a\in\Gamma^\tau\Longleftrightarrow\exists n,\ \pi_n(a)=e\Longleftrightarrow$ il existe un sous-
\deleted{tel que $a$}
\deleted{(ii) Soit} groupe \deleted{ouvert} [néc. ouvert car $\Gamma$ sép sur $k$]
$\Gamma'$ \uline{fini} sur $k$ tel que $a\in\Gamma'\Longleftrightarrow$
l'un des $a^n$ est \uline{fermé}.

(iii) Si $\Gamma$ est commutatif, alors $\Gamma^\tau$ est le
sup des sous-groupes [ouverts] \uline{finis}
de $\Gamma$. C'est aussi le plus petit sous-groupe
tel que $\Gamma/\Gamma^\tau$ soit ``sans torsion''

(iv) Soit $G^\tau$ l'image inverse de $\Gamma^\tau$. \deleted{\unsure{Alors}}\inserted{C'est}
un sous-schéma ouvert de $G$, et
$a\in G^\tau\Longleftrightarrow\exists n,\ \pi_n(a)\in G_0\Longleftrightarrow$ Il existe
un sous-groupe ouvert $H$ de $G$ contenant \unsure{$a$|les $a$} \deleted{\unsure{$a^n$|les}}
$\Longleftrightarrow\exists$ \unsure{ss-ensemble|ss-schéma} \uline{de type fini} de $G$ contenant les $a^n$
% archive sheet number: 205

\scanpage{208}
\uline{Critère de type fini}
\margin{$\deleted{\unsure{\beta}}\unsure{\mu|\beta}(s)+2\alpha(s)$ \unsure{$1/2$} continue}

\uuline{Proposition} \quad Soit $f:X\longrightarrow Y$ un morphisme loc. de
type fini de préschémas loc noethériens. On
suppose (i) $f$ séparé (ii) $f$ \deleted{\unsure{univers.}} ouvert (iii) les
fibres de $f$ sont \deleted{\unsure{de type fini}}\inserted{propres sur les $k(y)$} Alors $f$
est de type fini.
% slanted
\margin{prouver que $\underline{\mathrm{Pic}}^{\unsure{\tau}}_{X/S}$ est quasi-compact si $f:X\unsure{-|\to}S$ est \unsure{comment|comme ça} ?}

\inserted{On peut supposer $Y$ affine, donc \unsure{noeth.}} Il suffit de prouver $f$ qu. cpct, i.e $X$ quasi cpct.
Une \unsure{récurrence|récurrance} noeth. \deleted{\unsure{nous}} \unsure{$1/$} nous
ramène à prouver ceci : \deleted{il existe si $Y$ \unsure{en}}\inserted{Il existe}
un ouvert \unsure{non} vide de $Y$ tel que $X|U$ soit
qu. cpct. \inserted{On peut supposer $Y$ intègre, pt gén. $\eta$} Il existe un ouvert qu. cpct $V$ de $X$ \deleted{de
type fini} contenant $X_\eta$. Il existe un
$X'$ \deleted{de type fini}\inserted{\uline{propre}} sur $Y$, et un \unsure{$k(\eta)$-}isom
$X'_\eta\simeq X_\eta=V_\eta$. Si on restreint $Y$, on peut
supposer que le dernier isom provient
d'un isomorphis $X'\simeq V$. Donc \deleted{$V$ \unsure{soit} propre}
$V$ est propre sur $Y$, donc fermé dans $X$ [\unsure{car|où} $X$ est
séparé sur $\deleted{\unsure{X}}Y$]. \deleted{\unsure{On a un pt}} \deleted{Alors $X$\unsure{$-V$|$\,{}_\eta$} est une partie}
\deleted{\unsure{f} donc} $X-V$ est \uline{ouvert}. \deleted{Comme $X$}\inserted{Comme $X\unsure{-|\to}Y$}
\deleted{\unsure{la} \unsure{lettre|loi} dernière \unsure{donne}} est ouvert, \deleted{$X-V$ donc}
si on avait $X-V\ne\varnothing$, $X-V$ dominerait $Y$,
\deleted{donc} ce qui est incompatible avec $X_\eta=V_\eta$
i.e. $(X-V)_\eta=\varnothing$. Donc $X=V$, d'où la conclusion
% archive sheet number: 207

\scanpage{209}
\uline{Corollaire} Soit $\deleted{X}Z$ une partie de $X$. On
suppose (i) $\deleted{\unsure{X}}Z$ loc. constructible dans $X$
(ii) $f:X-S$ séparé loc. de type fini (iii) \deleted{$\unsure{X|Z}\to S$}
\inserted{Pour tt $x\in Z$, et $y'$ gén de $f(x)$, $\exists$ gén $x'\in Z$ de $x$, au dessus de $y'$}
\deleted{est ouvert en les $x'\in\unsure{X|Z}$} (iv) les fibres de
$\deleted{\unsure{X}}Z$ sur $S$ sont \uline{propres}. Alors $\deleted{\unsure{X}}Z-S$ est
qu. cpct.

On pt supposer $S$ affine \unsure{donc|ou} noeth, et on est
ramené à prouver que \inserted{si $S\ne\varnothing$,} il existe un ouvert $U\ne\varnothing$
de $S$ tel que $\deleted{\unsure{X}}Z\cap f^{-1}U$ soit qu. cpct. On peut supposer
\unsure{$S$ irréduct. de pt gén $\eta$.}

On construit un $Z'$ de type fini sur $S$ et un
isom $Z'_\eta\to Z_\eta$, \deleted{on peut supposer $Z'$, propre} \inserted{il existe un ouvert qu. cpct}
$V$ de type fini de $X$ contenant $Z_\eta$. Restreignant $S$,
on peut supposer $Z'_\eta\to Z_\eta\subset V$ provient d'un
morphisme $Z'\to\deleted{\unsure{X|Z}}V$, [que ce dernier est
une immersion], \inserted{que $Z\cap V=Z'$} que $Z'$ est \uline{propre} sur $S$.

Alors $Z'$ est \uline{fermé} dans $X$, et par suite on
aura $Z=(Z\cap V=Z')\cup(Z-Z')$, décomposition
\inserted{de $Z$} en deux \uline{ouverts \inserted{relatifs} disjoints}. On est
ainsi ramené au cas où $Z_\eta=\varnothing$,
qui implique $Z=\varnothing$ ....
% archive sheet number: 208
